\documentclass[12pt,a4paper]{article}
\usepackage[italian]{babel}
\usepackage[utf8]{inputenc}
\usepackage[T1]{fontenc}
\usepackage{geometry}
\geometry{margin=2.5cm}
\usepackage{booktabs}
\usepackage{array}
\usepackage{xcolor}
\usepackage{fancyhdr}
\usepackage{tikz}
\usepackage{colortbl}
\usepackage{hyperref}

\definecolor{cesareblu}{RGB}{30, 90, 160}
\definecolor{cesaregrigio}{RGB}{240, 240, 240}

\pagestyle{fancy}
\fancyhf{}
\rhead{Informatica --- Classe Prima/Seconda}
\lhead{Il Cifrario di Cesare}
\cfoot{\thepage}

\title{
  {\large Modulo: Introduzione alla Crittografia}\\[0.5em]
  {\LARGE \textbf{Il Cifrario di Cesare}}\\[0.5em]
  {\large Messaggi segreti dai Romani a oggi}
}
\author{Dispensa per gli studenti}
\date{}

\begin{document}
\maketitle
\thispagestyle{fancy}

% -------------------------------------------
\section*{Perché la crittografia?}

Nel film \textit{The Imitation Game}, i nazisti usavano una macchina
chiamata \textbf{Enigma} per cifrare i messaggi militari. Alan Turing
doveva capire il codice senza conoscerlo.

La domanda che ci facciamo oggi è più semplice, ma stessa idea:

\begin{center}
  \large\textbf{Come mando un messaggio segreto che solo il mio amico
  può leggere?}
\end{center}

% -------------------------------------------
\section{Il Cifrario di Cesare}

Giulio Cesare mandava ordini segreti ai suoi generali usando un
trucco semplicissimo: \textbf{spostare le lettere dell'alfabeto}.

\subsection{La tabella di cifratura (chiave = 3)}

Per cifrare un messaggio con chiave 3, scrivi l'alfabeto normale
nella riga in alto. Sotto, scrivi lo stesso alfabeto ma partendo
dalla lettera D (spostato di 3):

\vspace{0.5em}
\begin{center}
\renewcommand{\arraystretch}{1.4}
\begin{tabular}{|c|c|c|c|c|c|c|c|c|c|c|c|c|c|}
\hline
\rowcolor{cesareblu!80}
\color{white}\textbf{Normale} &
\color{white}A & \color{white}B & \color{white}C &
\color{white}D & \color{white}E & \color{white}F &
\color{white}G & \color{white}H & \color{white}I &
\color{white}L & \color{white}M & \color{white}N &
\color{white}O \\
\hline
\rowcolor{cesaregrigio}
\textbf{Cifrato} &
D & E & F & G & H & I & L & M & N & O & P & Q & R \\
\hline
\end{tabular}

\vspace{0.3em}

\begin{tabular}{|c|c|c|c|c|c|c|c|c|c|c|c|c|}
\hline
\rowcolor{cesareblu!80}
\color{white}\textbf{Normale} &
\color{white}P & \color{white}Q & \color{white}R &
\color{white}S & \color{white}T & \color{white}U &
\color{white}V & \color{white}Z & \color{white}A &
\color{white}B & \color{white}C & \color{white}D\\
\hline
\rowcolor{cesaregrigio}
\textbf{Cifrato} &
S & T & U & V & Z & A & B & C & D & E & F & G \\
\hline
\end{tabular}
\end{center}
\vspace{0.5em}

\textbf{Per cifrare:} guarda la lettera nella riga normale $\rightarrow$
scrivi quella nella riga cifrata.\\
\textbf{Per decifrare:} guarda la lettera nella riga cifrata $\rightarrow$
risali a quella nella riga normale.

\subsection{Esempio pratico}

\vspace{0.5em}
\begin{center}
\renewcommand{\arraystretch}{1.6}
\begin{tabular}{lccccccc}
\toprule
\textbf{Testo normale} & C & I & A & O & \quad & M & O \\
\textbf{Testo cifrato} & F & L & D & R & \quad & P & R \\
\bottomrule
\end{tabular}
\end{center}
\vspace{0.5em}

Il messaggio \texttt{CIAO MO} diventa \texttt{FLDR PR}.
Chi non conosce la chiave non capisce niente!

% -------------------------------------------
\section{Perché è facile da rompere?}

La chiave è solo un numero da 1 a 25. Un attaccante può
semplicemente provarle tutte, una per una, finché il messaggio
non torna a avere senso. Questo si chiama \textbf{attacco a forza bruta}.

Con carta e penna ci vogliono pochi minuti. Con un computer, meno
di un secondo. Ecco perché oggi usiamo sistemi molto più complicati!

% -------------------------------------------
\section{Attività: Agenti Segreti}

\subsection*{Materiale}
\begin{itemize}
  \item Un messaggio cifrato per ogni squadra (stesso per tutti)
  \item Un foglio con la tabella alfabeto vuota (da compilare)
  \item Una busta chiusa con la chiave segreta
\end{itemize}

\subsection*{Regole}
\begin{enumerate}
  \item Il professore sceglie una chiave e la chiude in una busta.
  \item La classe si divide in squadre da 3--4 studenti.
  \item Ogni squadra prova le chiavi una alla volta, compilando
        la tabella ogni volta.
  \item La prima squadra che decifra il messaggio e indovina
        la chiave vince.
  \item Si apre la busta per confermare.
\end{enumerate}

\subsection*{Trucco}
Se il testo decifrato ha senso in italiano, hai trovato la chiave!

\vspace{1em}
\begin{center}
\renewcommand{\arraystretch}{1.8}
\begin{tabular}{|c|*{14}{c|}}
\hline
\textbf{Normale} &
A&B&C&D&E&F&G&H&I&L&M&N&O\\
\hline
\textbf{Cifrato} &
\phantom{A}&\phantom{A}&\phantom{A}&\phantom{A}&\phantom{A}&
\phantom{A}&\phantom{A}&\phantom{A}&\phantom{A}&\phantom{A}&
\phantom{A}&\phantom{A}&\phantom{A}\\
\hline
\end{tabular}

\vspace{0.4em}

\begin{tabular}{|c|*{13}{c|}}
\hline
\textbf{Normale} &
P&Q&R&S&T&U&V&Z&A&B&C&D\\
\hline
\textbf{Cifrato} &
\phantom{A}&\phantom{A}&\phantom{A}&\phantom{A}&\phantom{A}&
\phantom{A}&\phantom{A}&\phantom{A}&\phantom{A}&\phantom{A}&
\phantom{A}&\phantom{A}\\
\hline
\end{tabular}
\end{center}

\vspace{1em}
\noindent\textbf{Messaggio da decifrare:}
\vspace{0.5em}

\begin{center}
\large\texttt{\_\_\_\_ \_\_\_\_ \_\_\_\_ \_\_\_\_ \_\_\_\_}
\end{center}
\textit{(il professore inserisce qui il messaggio cifrato)}

\end{document}